\chapter{マシンの入力：目、耳、その他のセンサはどう接続されているか}
\chaptermark{マシンの入力}
\label{sec:sensors}

\section{変換器としてのセンサ}

図~\ref{fig:machine}では、データは左側の「センサ」ブロックからマシンに入ってきた。今度はそのブロックを開けてみよう。エンジニアから見れば、どの感覚器官も\emph{変換器}である。アナログの入力段があり、最後にA/D変換器がある。ただし出力のデジタル値は数ではなく、スパイクである。

感覚器官の構成はどれも同じである（図~\ref{fig:transducer}）。光、圧力、振動、分子といった物理量が、感覚細胞の膜にある受容体タンパク質に作用する。このタンパク質はゲートとして働き、自らイオンチャネルを開くか、短い反応の連鎖を介して開く。\emph{受容器電流}が流れ、それとともに膜にアナログ電圧が生じる。電圧はゲイン調整を経てスパイク符号器に入る。符号器はおなじみの積分型ニューロン~\eqref{eq:glutneuron}で、レベルを発火率とスパイクの時刻に変える。出力はほぼ例外なく同じである。目、耳、嗅覚、皮膚の感覚ニューロンはグルタミン酸を放出するので、入力はそのまま\nt{GLUT}バスに接続される。

\begin{figure}[t]
\centering
\begin{tikzpicture}[>=Stealth,
  blk/.style={draw,very thick,rounded corners=2pt,align=center,font=\footnotesize,minimum height=1.0cm,text width=1.9cm,fill=blue!5}]
\node[align=center,font=\small] (P) at (0.1,0) {光、音、\\圧力、\\分子};
\node[blk] (R) at (3.0,0) {受容体\\ゲート};
\node[blk] (G) at (6.5,0) {ゲイン\\調整};
\node[blk] (E) at (10.0,0) {スパイク\\符号器};
\node[align=center,font=\small] (B) at (13.4,0) {\nt{GLUT}\\バス};
\draw[->,very thick] (P) -- (R);
\foreach \ic/\xx in {sun/-1.1,speaker/-0.3,press/0.5,molecule/1.3}{\pic[scale=0.85] at (\xx,1.05) {\ic};}
\draw[->,very thick] (R) -- node[above,font=\scriptsize] {電流} (G);
\draw[->,very thick] (G) -- node[above,font=\scriptsize] {レベル} (E);
\draw[->,very thick,red!70!black] (E) -- node[above,font=\scriptsize,black] {スパイク} (B);
\draw[decorate,decoration={brace,amplitude=5pt,mirror},gray!60!black] (1.8,-0.8) -- (7.7,-0.8) node[midway,below=5pt,font=\scriptsize,black] {アナログ部};
\draw[decorate,decoration={brace,amplitude=5pt,mirror},gray!60!black] (8.8,-0.8) -- (11.2,-0.8) node[midway,below=5pt,font=\scriptsize,black] {スパイク};
\end{tikzpicture}
\caption{変換の連鎖としての感覚器官。受容体タンパク質がチャネルを開いて電流を生む。ゲイン調整がそれを状況に合わせる。符号器~\eqref{eq:glutneuron}がレベルをスパイクに変え、スパイクは\nt{GLUT}バスへ出ていく。目と耳では最初の段はスパイクをまったく出さない。遠くへ送る必要が生じるまで、信号はアナログのままである。}
\label{fig:transducer}
\end{figure}

ある細部は電子回路の技術者にはなじみ深い。目と耳では、いちばん最初の細胞はスパイクをまったく出さない。基板上のアナログ段のように、なめらかな電圧を隣の細胞に渡すだけである。スパイクが現れるのは、信号を遠くへ運ぶ必要がある場所だけである。アナログ信号は数ミリの距離で減衰し、ノイズを拾う。一方スパイクは、第\ref{sec:neurontypes}章で見たように、同じ高さのまま配線の終端まで届く。デジタル化は長い線路が始まる場所で行われる。

\section{物理の限界で}

脳のセンサは感度の物理的限界で働いている。暗所の光センサは\emph{光子1個}に応答する。光の粒子1個の吸収が約1ピコアンペアの電流を生み、それは自身のノイズを背景にしても検出できる~\cite{Baylor1979}。音のセンサは、自分の感覚毛束の1ナノメートル未満の変位を検出する~\cite{Hudspeth1989}。

光が少ないとき、精度を制限するのはセンサではなく光そのものである。光子はポアソン流としてランダムに到来する。蓄積時間 $T$ の間にセンサに平均 $N=\Phi T$ 個の光子が入るなら、その数は $\sqrt N$ だけゆらぐ。増分 $\Delta N$ を検出するにはこのゆらぎの数倍でなければならず、識別できる明るさの変化の割合は次のようになる：
\begin{empheq}[box=\keybox]{equation}
\frac{\Delta I}{I} \;\approx\; \frac{k}{\sqrt{\Phi\,T}} .
\label{eq:photonnoise}
\end{empheq}
これはスパイク数についての式~\eqref{eq:rateerror}と同じである。光子のショットノイズとスパイクのショットノイズは同じ法則に従う。暗所で目が精度を上げる方法は一つしかない。光子をより長く蓄積することであり、実際に蓄積時間は十分の数秒まで伸びる。だから薄暗がりではものの見え方が遅くなる。

\section{ダイナミックレンジ：ゲイン調整}

センサにとって最も難しい工学的課題はレンジである。照度は星明かりの夜から日の当たる雪原まで約10桁変わり、音の強さは聴覚のしきい値から痛みの限界まで12桁変わる。一方、スパイクの発火率はゼロから毎秒数百までで、3桁にも満たない。このような入力をこのような出力に線形に収めることはできない。

解決法はどの受信機とも同じ、\emph{自動ゲイン調整}である。目の感覚細胞の応答は次の式でよく記述される：
\begin{empheq}[box=\keybox]{equation}
r \;=\; r_{\max}\,\frac{I}{I+\sigma} ,
\label{eq:nakarushton}
\end{empheq}
ここで $\sigma$ は応答が半分に達する明るさである~\cite{NakaRushton1966}。\eqref{eq:nakarushton}自体は2〜3桁しかカバーしない。だが $\sigma$ は一定ではない。明るい背景で長く働くと、平均の明るさを追って大きくなる。応答曲線は明るさの対数軸に沿ってずれ、そのつど急峻な部分を現在の入力の位置に置く（図~\ref{fig:adaptation}）。これは第\ref{sec:gaba}章のシャント抑制と同じ、全体の活動による除算である~\cite{Carandini2012}。ただしここでの分母は、直近数秒の平均の明るさである。

\begin{figure}[t]
\centering
\begin{tikzpicture}[>=Stealth,x=1.25cm,y=2.8cm]
\draw[->,gray!70!black] (-2,0) -- (6.4,0) node[below left,font=\small,black] {$\log_{10} I$};
\draw[->,gray!70!black] (-2,0) -- (-2,1.15) node[left,font=\small,black] {$r/r_{\max}$};
\foreach \u in {-2,0,2,4,6}{\draw[gray!60!black] (\u,-0.02) -- (\u,0.02); \node[below,font=\scriptsize] at (\u,0) {$\u$};}
\draw[densely dashed,gray!60!black] (-2,0.5) -- (6.2,0.5);
\foreach \s/\c/\lab in {0/blue!60!black/{暗い背景},2/green!45!black/{中程度},4/red!70!black/{明るい}}{
  \pgfmathsetmacro{\lo}{max(-2,\s-3)}
  \draw[very thick,\c] (-2,0) -- (\lo,0) plot[domain=\lo:6.2,samples=80] (\x,{1/(1+10^(\s-\x))});
  \node[font=\scriptsize,\c,anchor=west] at (\s+0.15,0.42) {\lab};}
\foreach \s/\ic in {0/moon,2/cloud,4/sun}{\pic at (\s+0.6,0.64) {\ic};}
\end{tikzpicture}
\caption{\eqref{eq:nakarushton}による光センサのゲイン調整。各曲線は明るさの2〜3桁をカバーする。より明るい背景に移ると $\sigma$ が大きくなり、曲線は対数軸に沿ってずれる。ずれていく曲線全体でレンジ全体をカバーする。}
\label{fig:adaptation}
\end{figure}

ゲイン調整には、第\ref{sec:code}章でおなじみの代償がある。センサが伝えるのは明るさではなく、\emph{背景に対する}明るさである。一定の入力は次第に応答を起こさなくなり、変化のたびに応答が起きる。マシンの入力は最初から、レベルではなく変化について語る。ここにも命題~\ref{prop:economy}と同じスパイクの節約がある~\cite{Barlow1961}。第\ref{sec:glutcomputer}章のオン型とオフ型のセンサもここから来る~\cite{Schiller1992}。一方は明るくなったことを、他方は暗くなったことを伝える。

エンジニアはこの手法を\emph{イベントカメラ}で再現した。このカメラはクロックに合わせてフレームを撮るのではない。各画素が、共通のクロックなしに独立して、明るさの対数が所定の割合だけ変化したときに「明るく」または「暗く」のイベントを出す~\cite{Lichtsteiner2008}。これはまさに第\ref{sec:glutcomputer}章のオン／オフの2線式符号をシリコンで実現したものであり、通常のイメージセンサには得られないレンジと速さをカメラに与える~\cite{Gallego2022}。

\section{目：1本のケーブルへの圧縮}

目にはおよそ $10^8$ 個の光センサがあるが~\cite{Curcio1990}、画像を脳へ運ぶ出力ニューロンは約100万個しかない。信号は目を出る前に約100分の1に圧縮される。圧縮の主な手法はすでにおなじみである。各出力ニューロンは、小さな斑点の明るさとその周囲の明るさとの差に応答する。これは第\ref{sec:gaba}章と同じ、抑制による近傍の減算である。画像の一様な部分はほとんどコストがかからず、エッジと変化が伝えられる。出力ニューロンは種類ごとに特化している。明るくなったことを伝えるもの、暗くなったことを伝えるもの、特定の方向の運動を伝えるものがある（図~\ref{fig:direction}）。マウスではこうした出力チャネルが30種類以上数えられている~\cite{Baden2016}。

できあがったケーブルの伝送容量はどれほどか。モルモットでは出力ニューロンの記録から、$\sim10^5$ 個のニューロンで毎秒約 $875$ キロビットが測定された。ヒトの100万個の出力ニューロンに換算すると、毎秒 $10^7$ ビットのオーダーになる~\cite{Koch2006}。昔の10メガビットのネットワークケーブルと同じである。ニューロンあたりの平均発火率が毎秒数スパイクなら、1スパイクあたり数ビットとなり、第\ref{sec:code}章で得た値と一致する。

\section{耳：フィルタバンク}

耳が解く課題は別である。音を周波数に分解しなければならない。エンジニアならバンドパスフィルタのバンクを置くだろう。耳はまさにそれを機械的に行う。音の振動は弾性のある隔壁の上を伝わり、隔壁の性質は長さに沿ってなめらかに変わるので、各部分がそれぞれの周波数で共振する。その部分にあるセンサは、自分の帯域にだけ応答する（図~\ref{fig:filterbank}）。周波数は場所に沿ってほぼ対数的に並ぶ。どのオクターブにも同程度の割合のセンサが割り当てられる。応答したセンサの番号が周波数であり、第\ref{sec:code}章の場所符号である。

\begin{figure}[t]
\centering
\begin{tikzpicture}[>=Stealth,x=1.35cm,y=2.2cm]
\draw[->,gray!70!black] (-0.3,0) -- (7.6,0) node[above left,font=\small,black] {周波数、Hz};
\draw[->,gray!70!black] (-0.3,0) -- (-0.3,1.15) node[left,font=\small,black] {応答};
\foreach \k/\f in {0/125,1/250,2/500,3/1000,4/2000,5/4000,6/8000}{
  \draw[gray!60!black] (\k,-0.02) -- (\k,0.02); \node[below,font=\scriptsize] at (\k,0) {\f};
  \draw[thick,blue!60!black] plot[domain={max(\k-1.2,-0.3)}:{\k+0.7},samples=60] (\x,{1/(1+((\x-\k)/(\x<\k?0.45:0.2))^4)});}
% a piano keyboard: seven white keys per octave, like the octaves on the axis
\foreach \i in {0,...,48}{\draw[gray!50!black,fill=white,line width=0.3pt] ({\i/7},-0.44) rectangle ({(\i+1)/7},-0.26);}
\foreach \o in {0,...,6}{\foreach \j in {0,1,3,4,5}{\fill[black] ({\o+(\j+1)/7-0.035},-0.35) rectangle ({\o+(\j+1)/7+0.035},-0.26);}}
\end{tikzpicture}
\caption{バンドパスフィルタのバンクとしての耳（模式図）。各曲線は、ある部分のセンサがさまざまな周波数の音にどう応答するかを示す。中心周波数は、対数目盛上の鍵盤のように1オクターブおきに並ぶ。実際のフィルタは高周波側の斜面が急で、低周波側がなだらかである。下は鍵盤で、耳と同じく、どのオクターブにも同じ幅が割り当てられている。}
\label{fig:filterbank}
\end{figure}

耳の共振器は受動的ではない。一部のセンサは自ら音に合わせて隔壁を揺らし、弱い振動を振幅で数千倍（$66$--$76$\,dB）増幅する。音が大きくなると増幅は下がる~\cite{Ruggero1997,Robles2001}。エンジニアから見れば、これは自励発振の境界ぎりぎりに置かれた正帰還増幅器である。第\ref{sec:gaba}章のネットワークと同じく、境界ぎわで生きている。境界の近くでは、系は小さな信号にとりわけ敏感になる。音量の圧縮も同じ増幅器が行う。小さな音は大きく、大きな音は小さく増幅される。

耳にはもう一つの符号がある。低い周波数では、ニューロンのスパイクは音と位相をそろえて出る。ニューロンは毎回の波で発火するわけではないが、発火するときは波の決まった部分で発火する。モルモットでは位相への同期は約 $600$\,Hz で弱まり始め、約 $3.5$\,kHz まではまだ認められる~\cite{PalmerRussell1986}。こうしてスパイクの時刻には音の正確な時間が保たれ、そこから両耳への音の到達時間差が得られる。第\ref{sec:glutcomputer}章のネットワークはこの差から方向を求める~\cite{Grothe2010}。高い周波数ではスパイクが波に追いつけず、場所符号だけが残る。一つの耳が第\ref{sec:code}章の二つの符号を両方使っている。ニューロンが追いつけるところでは時間を、追いつけないところでは場所を。

\section{その他のセンサ}

\begin{table}[t]
\centering
\footnotesize
\setlength{\tabcolsep}{3pt}
\renewcommand{\arraystretch}{1.15}
\begin{tabular}{>{\raggedright}p{1.9cm}>{\raggedright}p{2.6cm}>{\raggedright}p{4.0cm}>{\raggedright\arraybackslash}p{4.6cm}}
\toprule
感覚 & 何を測るか & 入力の仕組み & 本書の手法 \\
\midrule
視覚 & 光 & $\sim10^8$ 個のセンサ、$\sim100$ 分の1に圧縮、$\sim10^7$ bit/s & ゲイン調整、近傍の減算、2本の配線、運動の方向 \\
聴覚 & 空気の圧力 & 機械的フィルタバンク、能動的増幅器 & 場所符号と時間符号、遅延 \\
触覚 & 圧力、振動 & 速く順応するセンサと遅く順応するセンサ & 「ハイパスフィルタとローパスフィルタ」の対 \\
温度 & 熱 & 温センサと冷センサが別々 & 2線式符号 \\
嗅覚 & 分子 & 受容体は $\sim400$ 種類、各センサは1種類 & 組み合わせ符号：匂いは応答した種類の集合 \\
味覚 & 食物中の物質 & 五つの味質：甘味、苦味、塩味、酸味、うま味 & ラベル付き回線；一部の細胞はグルタミン酸ではなくATPやセロトニンを放出 \\
身体の位置 & 筋肉の伸張 & 伸張の長さと速さのセンサ & 運動制御器へのフィードバック \\
\bottomrule
\end{tabular}
\caption{センサはどう接続されているか。3列目の仕組みはすべて測定済みである。右の列はそれらを前の章の手法と結びつける。}
\label{tab:sensors}
\end{table}

残りの感覚は表~\ref{tab:sensors}にまとめた。ほとんどどの行にも、前の章の手法が見てとれる。触覚は一対のフィルタを使う。ある圧センサは接触に応答してすぐに沈黙し、別の圧センサは圧力がある間ずっと応答を保つ。温度は温と冷で別々の2本の配線で伝えられる。最も変わった入力は嗅覚である。ヒトの嗅覚受容体は約400種類あり、各嗅覚ニューロンはそのうち1種類だけを持つ~\cite{Mombaerts2004}。匂いは1種類ではなく一組の種類を活性化し、メッセージは\emph{どの}種類が応答したかである。エンジニアから見れば、これは長さ約400の2値ベクトルであり、とりうる値の数は天文学的である。

\begin{keyblock}
\begin{proposition}[マシンの入力]
\label{prop:sensors}
どの感覚も変換器であり、感度の物理的限界で働き、巨大なダイナミックレンジを自動ゲイン調整で圧縮し、\nt{GLUT}バスを通じて何よりも\emph{変化}を伝える。そのためにセンサはマシン自身と同じ手法を使う。2線式符号、場所符号、時間符号、そして背景による除算である。センサの仕組みは測定されている。その符号が1スパイクあたりの情報を最大にするよう調整されているというのは、十分な根拠のある仮説である。
\end{proposition}
\end{keyblock}

入力も、他のすべてと同じく非同期に働く。各センサは伝えることがあるときにスパイクを出し、感覚ごとに遅延も違う。音は数ミリ秒後に、光は数十ミリ秒後に届く。共通のクロックがないのに、マシンがそれらをどう一つの世界像にまとめるのか。これは第\ref{sec:synthesis}章の時間的協調の課題の一つである。

\section{問題}

\begin{exercise}[\lvlA]
\label{ex:11:1}
\emph{光子をどれだけ蓄積するか。} 薄暗がりでセンサに毎秒 $100$ 個の光子が入る。増分はゆらぎ~\eqref{eq:photonnoise}の $3$ 倍あれば検出できる。(a)~1個のセンサが明るさの $10\,\%$ の変化を検出するのにどれだけの時間が要るか。(b)~$0.2\,\text{s}$ で間に合わせるには何個のセンサを足し合わせればよいか。空間分解能は何分の1に落ちるか。
\end{exercise}

\begin{exercise}[\lvlA]
\label{ex:11:2}
\emph{1スパイクに何ビットあるか。} (a)~目は $10^6$ 個の出力ニューロンで、毎秒 $3$--$5$ スパイクの発火率で $\sim10^7$ bit/s を送る。$\Delta t=1\ms$ での容量~\eqref{eq:spikeentropy}のどれだけを使っているか。(b)~ある匂いが $\approx400$ 種類の受容体のうち $20$ 種類を活性化する。この組は何ビットを運ぶか。二つのランダムな匂いに共通する種類は平均で何個か。
\end{exercise}

\begin{exercise}[\lvlB]
\label{ex:11:3}
\emph{1本の曲線~\eqref{eq:nakarushton}で何ができるか。} (a)~応答が $r_{\max}$ の $10\,\%$ から $90\,\%$ まで上がるのは、明るさのどの範囲か。それは何桁か。(b)~明るさの10桁をカバーするには、ずれていく曲線の位置がいくつ要るか。音の大きさの12桁ではどうか。
\end{exercise}

\begin{exercise}[\lvlB]
\label{ex:11:4}
\emph{耳の時間符号の限界。} 両耳への音の到達時間差は最大 $0.22\,\text{m}/343\,\text{m/s}$ である。(a)~スパイクが純音と位相をそろえて出るとき、どの周波数まで方向を一意に決められるか。(b)~スパイクの時刻は $0.5\ms$ ゆらぐが、$20$~\textmu s を区別する必要がある。毎秒 $100$ スパイクのニューロンを $50\ms$ の間に何個平均すればよいか。
\end{exercise}

\begin{exercise}[\lvlB]
\label{ex:11:5}
\emph{目のケーブルに載せるイベントカメラ。} $10^6$ 画素、出力 $10^7$ bit/s のカメラが、画素の $\ln I$ が $\ln1.2$ だけ変化したときにイベント（アドレスと符号）を送る。長さ $500$ 画素、明るさの段差 $4$ 倍のエッジは、最大どれだけの速さで動けるか。同じ出力で、1画素 $8$ ビットの通常のカメラなら毎秒何フレームになるか。
\end{exercise}

\begin{exercise}[\lvlB]
\label{ex:11:6}
\emph{シミュレーション：ゲイン調整。} センサ~\eqref{eq:nakarushton}が $\sigma$ を $\tau=1\,\text{s}$ で、対数で $\tau\,\dd\ln\sigma/\dd t=\ln I-\ln\sigma$、または線形に $\tau\,\dd\sigma/\dd t=I-\sigma$ と調整する。背景は $1\to100\to1$ と跳ぶ。それぞれの法則で、上への跳びと下への跳びの後、$+10\,\%$ のプローブへの応答が順応後の値の半分まで回復するのに何秒かかるか。
\end{exercise}

\begin{exercise}[\lvlC]
\label{ex:11:7}
\emph{曲線はどんな世界に合わせてあるか。} 出力のノイズが小さく一定のとき、明るさについての情報が最大になるのは $r(I)=r_{\max}\Phi_I(I)$ のときである。ここで $\Phi_I$ は明るさの分布関数である。(a)~曲線~\eqref{eq:nakarushton}が最適となるのはどんな分布か。その $\log_{10}I$ の広がりはどれほどか。(b)~出力のノイズがポアソン型のとき、最適な曲線はどれか。
\end{exercise}
